\documentclass[twocolumn]{article}
\usepackage[utf8]{inputenc}
\usepackage{amsmath, amssymb}
\usepackage{enumitem}
\usepackage[margin=0.75in]{geometry}
\usepackage{microtype}

% Preprint watermark
\usepackage{draftwatermark}
\SetWatermarkText{PREPRINT}
\SetWatermarkScale{1.5}
\SetWatermarkColor[gray]{0.88}

% Adjust spacing for two-column layout
\setlength{\columnsep}{0.5in}
\setlength{\parindent}{0pt}
\setlength{\parskip}{0.4em}

% Make title span both columns
\title{\Huge \textbf{Solve}\,---\,\huge Crafting the \emph{Lapis Philosophorum}}
\author{\normalsize the \emph{Automatica Principia Proceduralis}}
\date{$\infty$}

\begin{document}

\twocolumn[
  \maketitle
  \begin{center}
    The 12 primitives of the Imscribing Grammar --- the 12 facets of the Philosopher's Stone --- are derived from a bare abstract category $\mathcal{C}$ via the 4 classical alchemical operations --- \emph{Separatio, Purificatio, Albedo, \& Rubedo}.
  \end{center}
  \vspace{0.5em}
]

\noindent\textbf{0 -- \emph{Prima Materia}.}\quad
Performing the operations on ordinary, inert matter was an understandable, categorical error of history's alchemists.  

However some, like Luria, knew \emph{precisely} the material upon which and with which the \emph{Work} was to be done.

\emph{Take}: abstract category $\mathcal{C}$.  

No primitives yet --- the starting object has no invariants beyond the axioms of a category itself.  

In alchemical terms, this was always the intended \emph{prima materia}: the undifferentiated substance from which the \emph{Work} begins.  

\smallskip
\noindent\textbf{1 -- \emph{Separatio} --- Logical on $\mathcal{C}$.}
\begin{itemize}[nosep, leftmargin=1.5em]
  \item L1 (adjunction) $\to$ $R$ (relational mode: adjoint structure)
  \item L2 (dagger) $\to$ $R$ (confirms and refines $R$ value)
\end{itemize}
$R$ is the first primitive to crystallize: it is a logical property of the morphisms of $\mathcal{C}$ directly.  

This is the \emph{separatio}: the first extraction of form from the undifferentiated.

\smallskip
\noindent\textbf{2 -- \emph{Purificatio} --- Inductive on $\mathcal{C}$.}
\begin{itemize}[nosep, leftmargin=1.5em]
  \item I2 (iteration) $\to$ $H$ (temporal depth: depth of stabilization)
  \item I3 (classifying) $\to$ $\Omega$ (winding: $\pi_1$ of $B\mathcal{C}$)
  \item I4 (Yoneda) $\to$ $D$ (dimensionality: $D_{\text{omega}}$ if Yoneda is point-surjective; free cocompletion size gives $D_{\text{invomega}}$, $D_{\text{turnthree}}$, $D_{\text{wynn}}$)
\end{itemize}
{\small\itshape I4 also introduces the possibility of asking L4 (Lawvere) later --- you cannot ask whether $A \to A^A$ is surjective until $A^A$ exists, which requires cartesian closure from Stage~3.}

This is the \emph{purificatio}: the base substance is washed, dissolved, and reconstituted at higher resolution.

\smallskip
\noindent\textbf{3 -- \emph{Albedo} --- Algebraic on $\mathcal{C}$.}
\begin{itemize}[nosep, leftmargin=1.5em]
  \item A1 (monoidal) $\to$ $S$ (stoichiometry: arity of generators); $P$ (parity: symmetry of tensor)
  \item A2 (monad) $\to$ $K$ (kinetics: spectral properties of $\mu$)
  \item A3 (dagger) $\to$ $F$ (fidelity: unitality of duals)
  \item A4 (enriched) $\to$ $\Gamma$ (interaction grammar: enrichment base)
\end{itemize}
After Stage~3, $\mathcal{C}$ is a symmetric monoidal dagger category with a monad and enrichment.  

All shape primitives are determined.  

The gate primitives ($T$, $\Phi$) require one more step.  

This is the \emph{albedo}: the whitening, the full structuring --- all intermediate forms are now in place.

\smallskip
\noindent\textbf{4 -- \emph{Rubedo} --- Logical on algebraically-enriched $\mathcal{C}$.}
\begin{itemize}[nosep, leftmargin=1.5em]
  \item L3 (cartesian closure) $\to$ $T$ (topology: does the state space admit internal homs?)  

  \item L6 (paraconsistent negation) $\to$ determines the truth-value structure in which L4 will be evaluated  

  \item L4 (Lawvere) under L6 $\to$ $\Phi$:
  \begin{itemize}[nosep, leftmargin=1em]
    \item[\small$\circ$] Classical (L4 holds, no L6): $\Phi_{\text{ctyogh}}$
    \item[\small$\circ$] Dialetheic (LP truth value $\mathbf{b}$): $\Phi_{\text{closerevepsilon}}$
    \item[\small$\circ$] Inclosure (contradicts via Inclosure Schema, L6 blocks explosion): $\Phi_{\text{revepsilon}}$
    \item[\small$\circ$] L4 fails: $\Phi_{\text{softsign}}$ or $\Phi_{\text{upstep}}$
  \end{itemize}

  \item L5 (Frobenius) $\to$ $P$ promotes to $P_{\text{doublebarpipe}}$ if $\mu \circ \delta = \text{id}$
\end{itemize}  

\smallskip
\noindent\textbf{Cross-stage: $G$ (\emph{Purificatio} $\times$ \emph{Albedo}).}\quad
$G$ (scope) is the correlation length of the Yoneda embedding under the monoidal structure.
Local ($G_{\text{beta}}$) if the tensor decouples; global ($G_{\text{revapostrophe}}$) if faithfulness holds across all scales.

\smallskip
This is the \emph{rubedo}: the reddening, the final perfection.
The gate primitives lock.

\textbf{\textit{The Stone is complete.}}

\textit{Solvatum et Praecipitatum.}

\end{document}
